% Part: computability
% Chapter: recursive-functions
% Section: partial-functions

\documentclass[../../../include/open-logic-section]{subfiles}

\begin{document}

\olfileid{cmp}{rec}{par}
\olsection{Fungsi Rekursif Parsial}


Kanggo menehi dhasar marang definisi fungsi rekursif,
gatekna yen bukti kita bab anane fungsi komputabel
kang dudu rekursif primitif sejatine mbuktekake luwih
akeh maneh. Argumene prasaja: kang kita gunakake
mung kasunyatan yen fungsi $f_0,f_1,\dots$ bisa
dienumerasi kanthi cara kang ndadekake, minangka
fungsi saka $x$ lan $y$, $f_x(y)$ komputabel.
Dadi argumen kasebut lumaku kanggo \emph{golongan
fungsi apa wae kang bisa dienumerasi kanthi cara
kaya mangkono}. Iki ndadekake kita nemoni masalah:
kita kepengin njlentrehake fungsi komputabel kanthi
eksplisit, nanging saben jlentrehan eksplisit ngenani
kumpulan fungsi komputabel ora bakal nyakup kabeh!

Carane metu saka masalah iki yaiku nglebokake
fungsi \emph{parsial}. Kita bakal weruh yen fungsi
komputabel parsial \emph{bisa} dienumerasi.
\iftag{TMs}{ Pancen, kita wis meh ngerti iki, amarga
  mesin Turing bisa dienumerasi kanthi sistematis.}{}
Kita bakal bali maneh marang argumen diagonal
lan nliti ngapa argumen kasebut ora bisa lumaku
nalika fungsi parsial uga dilebokake.

Saiki pitakonane: apa kang kudu ditambahake
marang fungsi rekursif primitif supaya kita
entuk kabeh fungsi rekursif parsial? Kita
kudu nindakake rong perkara:
\begin{enumerate}
\item Owahana definisi fungsi rekursif primitif
  supaya uga ngidini fungsi parsial.
\item \emph{Tambahna} sawenehing bab marang definisi
  supaya kalebu fungsi parsial anyar.
\end{enumerate}

Perkara kapisan gampang. Kaya sadurunge,
kita miwiti saka fungsi nol, penerus, lan
proyeksi, banjur nutup golongan iki marang
komposisi lan rekursi primitif. Bedane mung
definisi komposisi lan rekursi primitif kudu
diowahi supaya ngidini sawenehing term ing
definisi ora nduwé nilai. Yen $f$ lan $g$
fungsi parsial, kita nulis $f(x) \fdefined$
kanggo nyatakake yen $f$ ditegesi ing $x$,
yaiku $x$ kalebu domain $f$; lan
$f(x) \fundefined$ kanggo nyatakake kosok
baline, yaiku $f$ ora ditegesi ing~$x$.
Kita nggunakake $f(x) \simeq g(x)$ kanggo
nyatakake yen $f(x)$ lan $g(x)$ kalorone
ora ditegesi, utawa kalorone ditegesi lan
padha. Notasi iki uga digunakake kanggo
term kang luwih rumit. Kita netepake
konvensi yen $h$ lan $g_0$, \dots,~$g_k$
kabeh fungsi parsial, mula
\[
h(g_0(\vec x),\dots,g_k(\vec x))
\]
ditegesi yen lan mung yen saben $g_i$ ditegesi
ing $\vec x$, lan $h$ ditegesi ing
$g_0(\vec x)$, \dots,~$g_k(\vec x)$.
Kanthi pangerten iki, definisi komposisi
lan rekursi primitif kanggo fungsi parsial
padha kaya sadurunge, kajaba tandha
'$=$' kudu diganti nganggo '$\simeq$'.

Bab kang ditambahake marang definisi fungsi
rekursif primitif kanggo entuk fungsi parsial
yaiku \emph{operator panggolekan tanpa wates}.
Yen $f(x,\vec z)$ iku fungsi parsial sembarang
ing wilangan asli, tegese $\mu x \; f(x,\vec z)$
minangka
\begin{quote}
  $x$ paling cilik kang ndadekake
  $f(0,\vec z), f(1,\vec z), \dots, f(x,\vec
  z)$ kabeh ditegesi, lan $f(x,\vec z) = 0$,
  yen $x$ kaya mangkono ana
\end{quote}
kanthi pangerten yen $\mu x \; f(x,\vec z)$
ora ditegesi yen ora ana. Iki netepake
$\mu x \; f(x,\vec z)$ kanthi unik.

\begin{explain}
Gatekna yen definisi iki ora nyebut
\iftag{TMs}{ mesin Turing utawa}{} algoritma
utawa model komputasi tartamtu. Nanging
kaya komposisi lan rekursi primitif, ana
gagasan operasional bab komputasi ing
panggolekan tanpa wates. Kanggo fungsi
parsial kang komputabel, argumen kang
ora nduwé nilai cocog karo input kang
komputasine ora mandheg. Tata cara
ngitung $\mu x \; f(x,\vec z)$ yaiku:
itung $f(0,\vec z), f(1,\vec z),
f(2,\vec z)$ nganti ana nilai nol
kang diasilake. Nanging yen salah siji
komputasi ing antarane ora mandheg,
komputasi $\mu x \; f(x,\vec z)$
uga ora mandheg.
\end{explain}

Yen $R(x,\vec z)$ iku relasi sembarang,
$\mu x \; R(x,\vec z)$ ditegesi minangka
$\mu x \; (1 \tsub \Char{R}(x,\vec z))$.
Kanthi tembung liya, $\mu x \;
R(x,\vec z)$ ngasilake nilai $x$ paling
cilik kang nyukupi $R(x,\vec z)$.
Dadi yen $f(x,\vec z)$ fungsi total,
$\mu x \; f(x,\vec z)$ padha karo
$\mu x \; (f(x,\vec z) = 0)$.
Nanging gatekna yen definisi wiwitan
luwih umum, amarga ngidini $f(x,\vec z)$
ora ditegesi ing saben papan (dene
fungsi karakteristik saka relasi
mesthi total).

\begin{defn}
Himpunan \emph{fungsi rekursif parsial}
yaiku himpunan paling cilik saka fungsi
parsial ing wilangan asli menyang
wilangan asli (kanthi aritas warna-warna)
kang ngemot fungsi nol, penerus, lan
proyeksi, lan katutup tumrap komposisi,
rekursi primitif, lan panggolekan tanpa wates.
\end{defn}

Mesthi wae, sawenehing fungsi rekursif parsial
uga total, yaiku ditegesi kanggo saben argumen.

\begin{defn}
\ollabel{defn:recursive-fn}
Himpunan \emph{fungsi rekursif} yaiku
himpunan fungsi rekursif parsial kang total.
\end{defn}

Fungsi rekursif kadhang diarani
'rekursif total' kanggo nandheske yen
fungsi iki ditegesi ing saben papan.

\end{document}
